\section{Formato mínimo de registro por hallazgo}

Cada hallazgo debe registrarse con los campos de la Tabla~\ref{tab:formato-registro}. El registro separa el hecho observado de la interpretación y exige indicar la evidencia que permitió validar o descartar la hipótesis. El formato es obligatorio para los diez casos y se complementa con la ficha de cada caso.

\begin{table}[H]
	\caption{Formato mínimo de registro por hallazgo}
	\label{tab:formato-registro}
	\centering
	\footnotesize
	\renewcommand{\arraystretch}{1.2}
	\begin{tabularx}{\textwidth}{@{}>{\raggedright\arraybackslash}p{0.20\textwidth} Y@{}}
		\toprule
		\textbf{Campo} & \textbf{Qué registrar} \\
		\midrule
		Identificador & Caso, número de evidencia y nombre del analista. \\
		Fecha/hora & Marca temporal y zona horaria; sincronización conocida. \\
		Origen/destino & IP, MAC cuando sea pertinente, puertos y protocolo; anonimizar si aplica. \\
		Observación & Hecho directamente observado, sin conclusiones prematuras. \\
		Hipótesis & Explicación posible y alternativas benignas consideradas. \\
		Validación & PCAP o registro consultado, filtro, comandos de inspección y resultado. \\
		Severidad & Impacto, probabilidad y justificación. \\
		Acción & Medida propuesta, responsable, riesgo de falsos positivos y verificación. \\
		Integridad & Nombre de archivo, tamaño y hash SHA-256 de la evidencia original. \\
		\bottomrule
	\end{tabularx}
\end{table}

\noindent\textbf{Consolidado de hallazgos.} Use la Tabla~\ref{tab:consolidado} para mantener la trazabilidad entre los casos y el informe final. Los hashes y la cadena de custodia se registran en el Anexo~A.

\begin{table}[H]
	\caption{Consolidado de hallazgos de la Parte I}
	\label{tab:consolidado}
	\centering
	\footnotesize
	\setlength{\tabcolsep}{3pt}
	\renewcommand{\arraystretch}{1.3}
	\begin{tabularx}{\textwidth}{@{}c c >{\raggedright\arraybackslash}p{0.12\textwidth} >{\raggedright\arraybackslash}p{0.15\textwidth} Y c Y@{}}
		\toprule
		\textbf{No.} & \textbf{Caso} & \textbf{Fecha/hora} & \textbf{Origen/destino} & \textbf{Observación} & \textbf{Severidad} & \textbf{Acción} \\
		\midrule
		2 & Caso 2 & 2026-11-10 16:51 & \texttt{172.28.50.10} $\to$ \texttt{172.28.50.53} UDP/53 & 80 consultas TXT de 121 caracteres con 100 caracteres Base32; entropía 4,49; 0 alertas en la línea base y 82 en el túnel. & Media & Forzar el resolver interno, registrar consultas, limitar la tasa y alertar por longitud y entropía. \\ \hline
		6 & Caso 6 & 2026-11-10 16:52 & \texttt{172.28.50.20} (falsificado) $\to$ \texttt{172.28.50.53} UDP/53, 123 y 1900 & Amplificación DNS 34,7$\times$, NTP 5,1$\times$ y SSDP 10,4$\times$; 379.320~bytes recibidos por la víctima sin consultas propias. & Alta & Cerrar la recursión, deshabilitar monlist, restringir SSDP, aplicar RRL y filtrado de origen. \\ \hline
		& & & & & & \\
		& & & & & & \\
		& & & & & & \\
		& & & & & & \\
		& & & & & & \\
		& & & & & & \\
		& & & & & & \\
		& & & & & & \\
		\bottomrule
	\end{tabularx}
\end{table}

\begin{nota}{Regla de evidencia}
Una anomalía no es una prueba de ataque. Antes de asignar severidad, confirme la observación con al menos dos fuentes independientes (por ejemplo, PCAP y registro del servidor) y documente las explicaciones benignas que no se pudieron descartar. Si la evidencia es insuficiente, indíquelo como limitación. La conservación y el análisis de la evidencia siguen las recomendaciones de gestión de incidentes de \textcite{nist80061}.
\end{nota}
